% ======================================================================================
%  Appendix — progression-audit / TAE
% ======================================================================================

\section{Audit implementation and cohort details}
\label{app:setup}

\paragraph{Estimator and intervals.} Every AUC we compute is a patient-grouped, repeated five-fold out-of-fold
estimate from an $\ell_2$-regularized logistic model, with a patient cluster-bootstrap 95\% CI (patients, not eyes,
are the resampled unit) \citep{diciccio1996bootstrap}. Eyes are scored, but both eyes of a patient always share a
fold, so no eye of the same patient straddles the split. Tabular modalities are the pre-specified summaries GRAPE ships: the
visual field, the retinal-nerve-fiber-layer (RNFL) thickness from optical coherence tomography (OCT), and intraocular
pressure (IOP). The fundus embedding is a frozen RETFound-CFP \citep{zhou2023retfound}
representation reduced to 16 principal components. The device is controlled with a covariate-adjusted AUC
\citep{janes2009aroc}, the standard estimator for holding a covariate fixed when comparing discrimination: it conditions
the score on the acquisition device (binary TOPCON-vs-Canon), so a model cannot use the device as a feature.

\paragraph{The camera-only floor and the device probe.} The camera-only floor is the same logistic estimator trained
on a single binary device indicator (is-TOPCON) and no image content; it is training-free in the sense that it uses no
pixels and no learned representation. The device-decoding probe scores the frozen fundus embedding out-of-fold against
the device label; we report it on the full-fundus crop (AUC \num{0.98}) and on the optic-disc crop the downstream
papers use (\num{1.00}). We use three fundus crops throughout: the full field of view (``full''), a crop centered on the
optic disc (``disc'', the crop that decodes at \num{1.00} and the one the fundus probe uses), and a tighter disc
region-of-interest (``ROI''); Table~\ref{tab:ladder} reports all three. The camera-only floor carries its own patient cluster-bootstrap CI
(\num{0.56}~\ci{0.48}{0.64} / \num{0.60}~\ci{0.48}{0.71} / \num{0.58}~\ci{0.49}{0.66}), and every device-controlled
fundus read lies within it (\S\ref{sec:findings}).

\paragraph{The GPPM CFP-only number.} GPPM \citep{zheng2025gppm} reports two fundus-only disc region-of-interest (ROI) reads in its Table~3:
an independent evidential-deep-learning \emph{Baseline} (a separate architecture GPPM adopts, not its own branch) at
\num{0.577}/\num{0.571}/\num{0.595}, the value we use in the floor triangulation (F1), which sits at the camera floor on
all three endpoints. The second is GPPM's own class-balanced CFP branch (\emph{CFP\_M}) at \num{0.668}/\num{0.727}/\num{0.555},
whose higher PLR2/PLR3 are biased artifacts (zero sensitivity on PLR3 and MD, a biased AUC on 6--16 positives). Its headline \num{0.849} is a visual-field$+$CFP fusion under a combined criterion, not a CFP-only result, so
the CFP-only read we cite is the reliable one and it sits at the camera floor. This number is reported under GPPM's own
eye-level, device-uncontrolled, 196-eye protocol. Since an eye-level split and an uncontrolled device only inflate an
AUC, \num{0.58} is a generous upper bound on GPPM's baseline-fundus signal (a patient-grouped, device-controlled
evaluation could only lower it), so its coincidence with our camera floor is a conservative comparison,
not a like-for-like one.

\paragraph{Reproducing GPPM's two-field mechanism.} GPPM pairs a baseline visual-field test with a single follow-up
test and a fundus photograph (the paper leaves the fundus visit timing under-specified), down-weighting the fundus
branch to \num{0.01}, so the visual-field pair carries its number. We reproduce that mechanism directly on GRAPE's visual-field series; no fusion architecture is needed, since
the field decline is the signal. From five features of a two-field pair (the two mean-sensitivity levels, their
difference, the two-point rate, and the span between them) we predict the fixed PLR2/PLR3/MD label, patient-grouped.
Pairing the baseline with the last visit (GPPM's in-window structure) scores \num{0.83}/\num{0.86}/\num{0.86}. Pairing it
with the second visit, the same two tests but both early, sits at chance (\num{0.53}/\num{0.53}/\num{0.51}). The
intermediate pairs rise monotonically as the second test moves later (mean span \num{0.8}$\rightarrow$\num{2.5} years),
so the lift tracks how far the second test reaches into the label window, not the count of two. The same early pair
predicting a transparent MD-slope event on the disjoint later window (visits 3--end) cannot be distinguished from chance
(\num{0.31}~\ci{0.00}{0.82}, 5 held-out events), consistent with the single-field landmark (a fixed-time forecast). This reproduces GPPM's
two-field mechanism, not its evidential-fusion architecture. Its own CFP-only branch already sits at the floor
(above), so the fusion adds in-window functional change, not baseline prognosis (\texttt{grape\_two\_vft.json}). The
visual-field labels also inherit the large test--retest variability of perimetry, which worsens with field damage. This
adds noise to the trend fit and caps the AUC achievable by any prognostic model, over and above the small event counts.

\paragraph{Reproduction recipe (F2).} We fine-tuned ResNet-50, EfficientNet-B4, DenseNet-121, and ConvNeXt-Tiny, plus
a soft-vote ensemble, with the Fouda-style imbalance recipe (optic-disc ROI crop, focal loss $\alpha{=}0.75$,
$\gamma{=}2.0$, class-balanced sampling, last-30\% fine-tune, 120 balanced mini-batches of 32 [15 epochs $\times$ 8
batches]). The
\emph{published} protocol is a single random eye-level 80/20 split over five seeds (no patient grouping, no device
control, no cross-validation); the \emph{patient-grouped} protocol is a five-fold pooled out-of-fold (OOF) estimate with the camera-only
reference. Eye-level entries in Table~\ref{tab:collapse} are the median (and max) over seeds. The reversion is not
an under-training artifact: the \emph{same} recipe reaches the published range on the eye-level split, the camera floor
is training-free, and a frozen linear probe and a fine-tuned single-CFP model both sit at the same floor (F1), so the drop is
the split, not the training budget.

\paragraph{Compute.} All audit analyses (frozen embeddings, logistic probes, window-length sweep, landmark, bootstrap) run on CPU in
minutes; the F2 fine-tuning runs on a single consumer GPU. No result requires re-scoring to reproduce the figures:
every number is read from a committed data panel.

\section{The device confound in detail}
\label{app:device}

The camera is confounded with the label (Fig.~\ref{fig:device}a), so a device-only model scores
\num{0.561}/\num{0.600}/\num{0.580} with no pixels. The naive disc-crop fundus-embedding AUC is no higher than that
floor, and controlling the device leaves it at the floor: PLR3 drops from \num{0.57} to \num{0.52} and MD from \num{0.59}
to \num{0.55}, with a weak PLR2 residual (\num{0.59}, its interval including the floor). The full representation ladder
(Table~\ref{tab:ladder}) shows the same across eight representations and three crops. Because the acquisition-device field ships with
GRAPE, this confound is controllable, but no published model controls for it. Is the camera indexing genuinely
more-advanced eyes (in which case a fundus model reading severity would be doing real prognosis), or is it an
exposure-time artifact? We checked directly (\texttt{grape\_camera\_casemix.json}): baseline severity is matched across
cameras, TOPCON vs Canon median baseline VF mean sensitivity \num{22.98} vs \num{22.26}, RNFL \num{80} vs \num{81}, IOP
\num{16.5} vs \num{18.5} (all Mann--Whitney $p>0.05$), so the camera does not photograph sicker eyes. The one metric that
approaches a difference, IOP ($p{=}0.068$), runs higher in the \emph{less}-progressing camera (Canon), so it argues
against a case-mix explanation rather than for one. But TOPCON eyes are
followed markedly longer (median \num{3.12} vs \num{2.03} years, $p<0.001$), and a slope's significance is mechanically
easier to reach with longer follow-up. The camera-progression association is thus an exposure-time artifact, not baseline
case-mix and not retinal prognosis, which sharpens rather than weakens the point.

\section{Serial fundus does not forecast}
\label{app:serial}

The serial-image models that report \num{0.87}--\num{0.96} feed a fundus \emph{sequence}. On GRAPE's follow-up
photographs, even the in-window serial-fundus signal (whose input overlaps the label window) is at or below chance on all
three endpoints, naive or device-controlled, and the follow-up embedding identifies the acquisition device at AUC
\num{0.96} while 41\% of eyes switch camera between visits. A serial-fundus model on GRAPE is therefore reading a changing
device at least as much as a changing eye, and there is no prospective fundus forecast here. The published numbers come
from the evaluation protocol, which we reproduce directly below.

\paragraph{Direct reproduction of TSDF.} TSDF's released code (\texttt{github.com/zongzi13545329/MICCAI}) provides a direct
check of the mechanism. It trains a DSMIL (dual-stream multiple-instance learning) aggregator on the full per-eye visit sequence of precomputed
MAE (masked-autoencoder) features (in-window by construction: every follow-up visit is fed, so the input spans the label
window) under an eye-level
\texttt{KFold(shuffle, random\_state{=}42)} with unweighted
cross-entropy. It then reports the mean over folds of each fold's best-epoch test AUC, with the best epoch chosen on the
test fold itself, so epoch selection is not held out. Re-running that exact recipe on its own released features recovers \num{0.866}/\num{0.954}/\num{0.917}
(vs.\ published \num{0.866}/\num{0.956}/\num{0.917}; one PLR3 fold reaches \num{1.00} on two positives). Fixing the three
defects independently (a \texttt{GroupKFold} by the 144 patients, epoch selection on a patient-disjoint validation
split, and a pooled out-of-fold AUC with a patient cluster-bootstrap CI) reverts them to \num{0.43}~\ci{0.33}{0.52} /
\num{0.54}~\ci{0.37}{0.74} / \num{0.55}~\ci{0.42}{0.67} (this is the split-controlled read; controlling the device and
circularity together as well leaves it at \num{0.45}/\num{0.56}/\num{0.52}, Table~\ref{tab:auditgrid}). A ladder of fundus-only serial models (a first-two-image
structural-difference model $\rightarrow$ simple change features $\rightarrow$ pooled RETFound embeddings $\rightarrow$
a temporal attention head) reproduces the range as the eye-level max over seeds and sits at the patient-grouped floor at
every rung, so the reversion is a property of the protocol, not the model. This rung is a
model-class stand-in for DiffSight-Former, not its architecture: on GRAPE DiffSight takes the first two fundus images
per eye (its $\leq$\num{5}-visit pre-onset prospective forecasting is its separate experiment on the SIGF onset dataset,
not GRAPE) and encodes inter-visit change through vessel and optic-disc/cup difference maps, none of which is released.
With only two visits the input is below the three-visit slope threshold, so it is non-circular but not prospective. We reproduce
the part that can be tested from the available materials: a small learned attention head over the \emph{standardized} frozen RETFound features of those
two images, letting the head model the inter-visit change rather than injecting a hand-made difference, with \emph{no}
hand-engineered features. It is trained in DiffSight's own regime (mini-batch, cosine annealing with warm restarts,
plain cross-entropy). Train-fold feature standardization conditions the head and lifts the whole seed distribution off
its earlier below-chance floor. Between the two protocols we change a single knob, the split, and hold the model,
input, and training fixed. Swept over eye-level seeds this head reaches all three reported numbers at the upper tail of
the seed distribution (PLR2/PLR3/MD max \num{0.82}/\num{0.97}/\num{0.90} vs.\ its \num{0.79}/\num{0.88}/\num{0.80}, on
\num{5}/\num{2}/\num{3} test positives; the reported values sit in the top decile of the seed distribution),
reproducing DiffSight's single-split numbers as the tail of a leaky regime. Changing only the split to patient-grouped,
the same head drops to the floor: \num{0.49}~\ci{0.38}{0.62} / \num{0.48}~\ci{0.23}{0.68} / \num{0.59}~\ci{0.43}{0.73},
every interval spanning chance and none clearing the camera floor. Because the reproduce and the revert differ in nothing but the split, the reported number is a property of the
split, not the image content. DiffSight's weights and its vessel/disc segmentation maps are unreleased, so this measures
its regime and model class rather than its exact architecture.

\paragraph{The baseline prospective forecast (well-powered).} The prescribed positive test predicts the shipped label
from the first visit, before the label's trend is computable, so it keeps the full event counts (\num{40}/\num{14}/\num{27})
and is well-powered, unlike a strictly-future landmark. We run it two ways, patient-grouped and device-controlled.
\emph{Image-only} (fundus alone): the first visit is the single-CFP disc probe (\num{0.60}/\num{0.57}/\num{0.59}), the
first two images the DiffSight two-visit head (\num{0.49}/\num{0.48}/\num{0.59}); both sit at the floor.
\emph{Fusion} (visual field, OCT-RNFL, fundus morphometry, and IOP): the first visit scores
\num{0.58}/\num{0.56}/\num{0.59}, every \num{95}\% interval spanning chance, and extending the input to the first two
visits stays at the floor.
This also completes the serial-model audit non-circularly: restricting each serial model to its first $\leq$2 visits is
the image-only read (TSDF, DiffSight) or the fusion read (GPPM's baseline$+$early visual-field$+$fundus), all at the
floor. A fundus-specific TSDF cannot be run more prospectively than this on GRAPE: its fundus follow-up is too sparse
($\approx$1 pre-window image per eye), so TSDF's high number intrinsically requires the in-window sequence. The
first-visit read is cleanly non-circular; the first-two read adds an early-trajectory point that stays below the
three-visit significance threshold.

\section{Dial robustness: the onset tracks the label, not the visit count}
\label{app:dialrobust}

If the dial's jump were an artifact of ``more visits, more signal'', it would sit at the same $k$ for every label. It
does not: it tracks when the label's defining statistic first becomes computable. Recomputing the dial with the same
window-feature, patient-grouped out-of-fold pipeline for alternative labels, the onset moves with the definition. On the
MD-slope label the jump is at $k{=}3$ under the shipped $p<0.05$: the AUC rises from
$0.47$~\ci{0.38}{0.55} at $k{=}2$ to $0.78$~\ci{0.67}{0.88} at $k{=}3$, disjoint intervals, and the jump stays at $k{=}3$
under a looser $p<0.10$ ($0.66$~\ci{0.57}{0.74} to $0.86$~\ci{0.80}{0.93}). On a placebo label defined by a two-visit
change (computable at $k{=}2$) the jump moves to $k{=}2$: $0.38$~\ci{0.28}{0.49} at $k{=}1$ to $1.00$~\ci{0.99}{1.00} at
$k{=}2$ (the window slope at $k{=}2$ is exactly the placebo's defining statistic, hence near-perfect recovery). The
intervals are disjoint across each onset, so it is predicted a priori by the label, which is the sense in which the dial
is a diagnostic rather than an observed curve (\texttt{grape\_dial\_robustness.json}).

\section{Full result tables}
\label{app:fulltables}

These tables give the exact numbers behind the figures. Table~\ref{tab:collapse} (reproduce-then-revert, single-CFP
architectures) backs Fig.~\ref{fig:revert}, and Table~\ref{tab:ladder} (the covariate-adjusted representation ladder)
backs the floor-as-ceiling claim (F1). Table~\ref{tab:auditgrid} gives the per-model audited AUCs with intervals
behind Fig.~\ref{fig:revert}, and Table~\ref{tab:dial} gives the circularity dial and per-statistic values behind
Fig.~\ref{fig:circular}.

% auto-generated by scripts/build_progression_audit_tables.py -- do not edit by hand
\begin{table}[h]
\caption{\textbf{Reproduce-then-revert (Check~2).} Fine-tuning the architectures the GRAPE/Fouda papers use, with their imbalance recipe, under the published eye-level protocol reaches the published range (the seed distribution's upper tail brackets the single published numbers), then reverts to the training-free camera-only floor under a patient-grouped split. Eye-level entries are median (max) over 5 seeds; grouped entries are patient-grouped five-fold pooled OOF. All values are AUC.}
\label{tab:collapse}
\centering
\small
\setlength{\tabcolsep}{5pt}
\begin{tabular}{lcccccc}
\toprule
 & \multicolumn{2}{c}{PLR2 (40 ev)} & \multicolumn{2}{c}{PLR3 (14 ev)} & \multicolumn{2}{c}{MD (27 ev)} \\
\cmidrule(lr){2-3}\cmidrule(lr){4-5}\cmidrule(lr){6-7}
Model & eye-level & grouped & eye-level & grouped & eye-level & grouped \\
\midrule
ResNet-50 & 0.41 (0.71) & 0.45 & 0.57 (0.72) & 0.48 & 0.51 (0.57) & 0.52 \\
EfficientNet-B4 & 0.57 (0.66) & 0.47 & 0.61 (0.83) & 0.44 & 0.45 (0.62) & 0.43 \\
DenseNet-121 & 0.53 (0.72) & 0.51 & 0.68 (0.83) & 0.56 & 0.70 (0.80) & 0.51 \\
ConvNeXt-Tiny & 0.43 (0.51) & 0.52 & 0.55 (0.71) & 0.46 & 0.41 (0.50) & 0.50 \\
Soft-vote ens. & 0.52 (0.69) & 0.47 & 0.67 (0.77) & 0.44 & 0.50 (0.53) & 0.49 \\
\midrule
\textit{published (lit.)} & \multicolumn{2}{c}{0.71} & \multicolumn{2}{c}{0.80 $\cdot$ 0.86 ens} & \multicolumn{2}{c}{0.73} \\
\textit{camera-only floor} & \multicolumn{2}{c}{0.56} & \multicolumn{2}{c}{0.60} & \multicolumn{2}{c}{0.58} \\
\bottomrule
\end{tabular}
\end{table}

\begin{table}[h]
\caption{\textbf{Representation ladder (Check~1).} Covariate-adjusted, patient-grouped repeated-OOF AUC [95\% CI] of frozen linear probes over eight representation/crop combinations, beside the training-free camera-only floor. Every representation's interval includes or sits below the floor on all three endpoints, except the lone ConvNeXt-B MD point whose lower bound clears it; that crossing does not replicate on ConvNeXt-L or under fine-tuning (\S\ref{sec:findings}).}
\label{tab:ladder}
\centering
\small
\setlength{\tabcolsep}{5pt}
\begin{tabular}{lccc}
\toprule
Representation (frozen LP, covariate-adjusted) & PLR2 & PLR3 & MD \\
\midrule
RETFound-CFP (full) & 0.504 \ci{0.39}{0.60} & 0.584 \ci{0.40}{0.74} & 0.558 \ci{0.44}{0.67} \\
RETFound-CFP (disc) & 0.588 \ci{0.49}{0.68} & 0.523 \ci{0.37}{0.68} & 0.552 \ci{0.44}{0.66} \\
RETFound-CFP (ROI) & 0.550 \ci{0.45}{0.65} & 0.646 \ci{0.51}{0.78} & 0.470 \ci{0.36}{0.59} \\
DINOv2-reg (full) & 0.286 \ci{0.20}{0.37} & 0.593 \ci{0.44}{0.74} & 0.480 \ci{0.35}{0.60} \\
DINOv2-reg (disc) & 0.480 \ci{0.39}{0.58} & 0.458 \ci{0.30}{0.62} & 0.556 \ci{0.41}{0.69} \\
ConvNeXt-T (disc) & 0.564 \ci{0.46}{0.66} & 0.626 \ci{0.47}{0.76} & 0.489 \ci{0.38}{0.60} \\
ConvNeXt-B (disc) & 0.593 \ci{0.51}{0.67} & 0.649 \ci{0.46}{0.81} & 0.744 \ci{0.66}{0.82} \\
ConvNeXt-L (disc) & 0.522 \ci{0.43}{0.60} & 0.523 \ci{0.40}{0.64} & 0.610 \ci{0.51}{0.71} \\
\midrule
\textit{camera-only floor} & 0.561 & 0.600 & 0.580 \\
\bottomrule
\end{tabular}
\end{table}

\begin{table}[h]
\caption{\textbf{Per-model audited AUCs [95\% CI] behind Fig.~\ref{fig:revert}.} Each reported number, reproduced by its own architecture, once each shortcut is controlled one at a time: the device (a covariate-adjusted AUC), a patient-grouped split, a non-circular input, and all three together. \emph{n/a}: the control does not apply, because a single- or two-image input has no serial window to leak; Fouda reports PLR3 only. Compare against the camera floor (\num{0.561}/\num{0.600}/\num{0.580}): no audited read clears it. GPPM's device and split reads are the highest of the five, but their intervals include the floor.}
\label{tab:auditgrid}
\centering
\small
\setlength{\tabcolsep}{5pt}
\begin{tabular}{ll ccc}
\toprule
Model & Control & PLR2 & PLR3 & MD \\
\midrule
Huang & device & 0.48 \ci{0.38}{0.59} & 0.49 \ci{0.30}{0.67} & 0.47 \ci{0.35}{0.60} \\
 & fair split & 0.44 \ci{0.34}{0.54} & 0.31 \ci{0.19}{0.45} & 0.54 \ci{0.42}{0.66} \\
 & circularity & \textit{n/a} & \textit{n/a} & \textit{n/a} \\
 & all three & 0.43 \ci{0.33}{0.53} & 0.27 \ci{0.15}{0.42} & 0.53 \ci{0.41}{0.66} \\
\addlinespace[2pt]
Fouda & device & -- & 0.43 \ci{0.23}{0.63} & -- \\
 & fair split & -- & 0.58 \ci{0.45}{0.71} & -- \\
 & circularity & -- & \textit{n/a} & -- \\
 & all three & -- & 0.54 \ci{0.41}{0.68} & -- \\
\addlinespace[2pt]
TSDF & device & 0.52 \ci{0.40}{0.62} & 0.57 \ci{0.42}{0.73} & 0.59 \ci{0.48}{0.68} \\
 & fair split & 0.43 \ci{0.33}{0.52} & 0.54 \ci{0.37}{0.74} & 0.55 \ci{0.42}{0.67} \\
 & circularity & 0.51 \ci{0.39}{0.62} & 0.60 \ci{0.45}{0.75} & 0.59 \ci{0.46}{0.71} \\
 & all three & 0.45 \ci{0.35}{0.55} & 0.56 \ci{0.44}{0.69} & 0.52 \ci{0.40}{0.64} \\
\addlinespace[2pt]
DiffSight & device & 0.52 \ci{0.37}{0.65} & 0.63 \ci{0.29}{0.88} & 0.48 \ci{0.36}{0.59} \\
 & fair split & 0.49 \ci{0.38}{0.62} & 0.48 \ci{0.23}{0.68} & 0.59 \ci{0.43}{0.73} \\
 & circularity & \textit{n/a} & \textit{n/a} & \textit{n/a} \\
 & all three & 0.49 \ci{0.37}{0.61} & 0.48 \ci{0.22}{0.69} & 0.56 \ci{0.42}{0.70} \\
\addlinespace[2pt]
GPPM & device & 0.67 \ci{0.55}{0.77} & 0.61 \ci{0.41}{0.87} & 0.52 \ci{0.38}{0.65} \\
 & fair split & 0.66 \ci{0.54}{0.77} & 0.57 \ci{0.31}{0.88} & 0.57 \ci{0.45}{0.68} \\
 & circularity & 0.56 \ci{0.44}{0.67} & 0.56 \ci{0.28}{0.79} & 0.48 \ci{0.33}{0.62} \\
 & all three & 0.55 \ci{0.43}{0.66} & 0.58 \ci{0.27}{0.83} & 0.52 \ci{0.38}{0.65} \\
\bottomrule
\end{tabular}
\end{table}

\begin{table}[h]
\caption{\textbf{The circularity dial behind Fig.~\ref{fig:circular}.} \emph{Top:} out-of-fold AUC for predicting the fixed label from the visual-field statistics of the first $k$ visits, climbing to \num{0.90} with the jump at $k{=}3$, when a slope and its $p$-value first become computable. \emph{Bottom:} full-series AUC from each single statistic against the MD label; the label-defining slope and its $p$-value carry the \num{0.90}.}
\label{tab:dial}
\centering
\small
\setlength{\tabcolsep}{6pt}
\begin{tabular}{lccccccc}
\toprule
$k$ input visits & 1 & 2 & 3 & 4 & 5 & 6 & full \\
\midrule
PLR2 & 0.58 & 0.51 & 0.68 & 0.73 & 0.82 & 0.82 & 0.82 \\
PLR3 & 0.56 & 0.55 & 0.74 & 0.77 & 0.85 & 0.86 & 0.87 \\
MD & 0.47 & 0.47 & 0.78 & 0.79 & 0.87 & 0.89 & 0.90 \\
\bottomrule
\end{tabular}
\vspace{4pt}
\centering
\begin{tabular}{lccccc}
\toprule
Single statistic & VF slope & slope $p$-value & last VF level & VF variability & \# visits \\
\midrule
MD-label AUC & 0.91 & 0.95 & 0.75 & 0.71 & 0.68 \\
\bottomrule
\end{tabular}
\end{table}



\section{A standardized re-evaluation protocol}
\label{app:template}

So that the audit is adoptable rather than only diagnostic, we release both a fixed evaluation split and a reporting
template. \emph{The split.} \texttt{grape\_standard\_splits.csv} assigns each of the 263 eyes to one of five folds by a
deterministic patient-grouped \texttt{GroupKFold} (\num{144} patients; both eyes of a patient in the same fold, so
fellow-eye leakage is impossible; one split serves PLR2/PLR3/MD). Scoring on these fixed folds removes the
single-favorable-split degree of freedom. \emph{The template.} On these folds, a trustworthy progression claim reports:
(i)~the training-free camera-only floor beside the model, and a device-adjusted AUC (we use a covariate-adjusted AUC);
(ii)~a patient-grouped, pooled out-of-fold AUC rather than a single split; (iii)~a baseline prospective forecast (the
first visit predicting the shipped label) alongside any in-window number; (iv)~a patient cluster-bootstrap \num{95}\% CI
and the event count for every endpoint; and (v)~whether the model input overlaps the label-defining window. Each item is
a one-line addition, and together they turn a leaderboard number into a defensible clinical claim.

\section{Reproducibility and broader impact}
\label{app:repro}

\paragraph{Reproducibility.} GRAPE is public under CC BY 4.0 (Figshare collection 6406319), and the RETFound-CFP weights
we probe are used non-commercially under CC BY-NC 4.0. On acceptance we release the analysis panel and the two
regeneration scripts under the MIT license. Every number in the paper is read from a
single committed data panel with a reconciliation gate, and the figures and tables are regenerated from it by two
scripts with no model calls or re-scoring; the audit estimators (device floor, device-decode probe, circularity sweep,
baseline prospective forecast, cluster-bootstrap CIs) are standard and specified in \S\ref{sec:audit} and \S\ref{app:setup}.
The panel is itself rebuilt from the committed per-model analysis outputs (small JSON, versioned alongside the code) by a
third script whose reconciliation gate cross-checks every locked constant; only regenerating those outputs from scratch
requires the frozen-feature tensors and the raw images on a GPU.

\paragraph{Broader impact.} The audit is a safeguard: it is designed to stop an over-optimistic progression-prediction
number from reaching a clinical decision. Its risk is the opposite: an over-strict reading could discourage a genuinely
useful model, so we frame the checks as necessary conditions for trusting a reported number, not as a ceiling on what
is clinically achievable, and we state explicitly that progression \emph{is} predictable at scale in cohorts without
these confounds. No patient-identifiable data are used beyond the public release.
